\documentclass[landscape]{tikzposter}
\geometry{paperwidth=36in,paperheight=48in}
\usepackage[english]{babel}
\usepackage{blindtext}
\usepackage{subfig}
\usepackage{listings}
\usepackage{adjustbox}
\usepackage{hyperref}
\title{\parbox{0.4\linewidth}{\centering Isosurface Extraction in a \\ Simple C/C++ Library}}
\institute{Idaho State University, College of Science and Engineering}
\author{Jonathan Glines, with undergraduate advisor Dr. John Edwards}

\makeatletter
\def\TP@titlegraphictotitledistance{-9cm}
\settitle{ \vbox{
\@titlegraphic \hfill \\ [\TP@titlegraphictotitledistance]
\centering
\color{titlefgcolor} {\bfseries \Huge \sc \@title \par}
\vspace*{1em}
{\huge \@author \par} \vspace*{1em} {\LARGE \@institute}
}}
\makeatother

\titlegraphic{
\includegraphics[trim={2.5in 3in 3in 2.5in},scale=1.4]{isu-wordmark.pdf}
 \hfill
}

\usetheme{Simple} % See Section 5
\definecolor{bengalOrange}{cmyk}{0,0.65,1,0}
\definecolorstyle{BengalColorStyle} {
\definecolor{colorOne}{named}{white}
\definecolor{colorTwo}{named}{bengalOrange}
\definecolor{colorThree}{named}{black}
}{
% Background Colors
\colorlet{backgroundcolor}{colorOne}
\colorlet{framecolor}{black}
% Title Colors
\colorlet{titlefgcolor}{black}
\colorlet{titlebgcolor}{colorTwo}
% Block Colors
\colorlet{blocktitlebgcolor}{colorThree}
\colorlet{blocktitlefgcolor}{colorTwo}
\colorlet{blockbodybgcolor}{white}
\colorlet{blockbodyfgcolor}{black}
% Innerblock Colors
\colorlet{innerblocktitlebgcolor}{white}
\colorlet{innerblocktitlefgcolor}{black}
\colorlet{innerblockbodybgcolor}{colorThree!30!white}
\colorlet{innerblockbodyfgcolor}{black}
% Note colors
\colorlet{notefgcolor}{black}
\colorlet{notebgcolor}{colorTwo!50!white}
\colorlet{noteframecolor}{colorTwo}
}
\usecolorstyle{BengalColorStyle}

\usetitlestyle{Filled}
\begin{document}
\maketitle % See Section 4.1
\begin{columns}
\column{0.33}
\block{\Huge Isosurface Extraction}{
\input{isosurfaceExtraction.tex}
}
\block{Original Marching Cubes}{
With the scalar field function $f$ in hand, we can readily generate a lattice
of sample values. The original marching cubes algorithm, along with most of its
variations, considers one voxel cube (i.e. eight samples arranged in a cube)
from this sample lattice at a time. From the eight sample values in this voxel,
we compute an 8-bit index according to the samples that are above or below the
isosurface.  The primary innovation of the marching cubes algorithm is the $2^8
= 256$ entry lookup table that enable us to quickly generate vertex positions
and triangles for our isosurface mesh.
}
\block{Accounting for Ambiguity}{
The original marching cubes algorithm as described by Lorensen and Cline uses
15 equivalence classes to generate the lookup tables. It was shown early in the
use of marching cubes that certain functions, such as those with saddle points,
could expose ambiguities in the way marching cubes was originally described.
More advanced treatments of the algorithm, such as Chernyaev's Marching Cubes
33 and Nielson's Asymptotic Decider, use additional equivalence classes that
resolve these ambiguities.
}
\block{Dual Marching Cubes}{
The dual of the marching cubes algorithm is easiest to understand by looking at
the popular voxel game Minecraft, since cube meshes generated by Minecraft are
topologically the same as any dual marching cubes mesh. The algorithm that
generates Minecraft-style meshes, called ``cuberille,"  does not perform any
interpolating or smoothing. More advanced dual methods will "relax" or "smooth"
the vertex positions of the cuberille mesh to produce a mesh much closer to the
actual isosurface.

%The advantages that dual methods have over ordinary marching cubes are
%immediately apparent: meshes generated by dual methods have more uniformly
%sized polygons than with ordinary marching cubes. Moreover, the dual methods
%generate quads, which are favored by 3D modelers over triangles.
}
\column{0.33}
\block{Transvoxel Algorithm}{
The Transvoxel isosurface extraction algorithm sets itself apart from ordinary
marching cubes by allowing seamless transitions between half and full
resolution samples. This algorithm was developed by Eric Lengyel in his PhD
dissertation titled ``Transition Cells for Dynamic Multiresolution Marching
Cubes."
}
\block{Cascading Transition Voxels}{
We are working to generalize the Transvoxel algorithm for transitioning between
$n$ levels of detail within the adaptive sample space of an octree. We hope to
apply this new algorithm to our research on generalized Voronoi diagram
computation.
}
\block{73 Transition Cell Cases}{
\newcommand{\transcellsize}{1in}
\def\transcellcolumns{
{"000", "00a", "010", "01a", "029", "03a", "04e", "061", "06b", "075", "0ab", "0ee", "15f"},
{"001", "00b", "011", "01b", "02a", "03b", "054", "062", "06c", "076", "0ad", "0ef"},
{"002", "00c", "012", "01c", "02b", "03d", "055", "063", "06e", "077", "0af", "0f5"},
{"003", "00d", "013", "01d", "02d", "044", "056", "065", "071", "07c", "0ba", "0fe"},
{"005", "00e", "015", "01e", "038", "045", "057", "066", "072", "07e", "0e5", "155"},
{"007", "00f", "017", "028", "039", "046", "05e", "06a", "073", "0aa", "0e7", "157"}}
\foreach \column in \transcellcolumns
{
  \begin{adjustbox}{valign=t}
  \begin{minipage}[t]{0.15\linewidth}
  \foreach \cell in \column
  {
    \begin{tikzfigure}
    \includegraphics[width=\transcellsize,height=\transcellsize]{../build/MC_TRANSVOXEL_transition_cell_0x\cell.png}
    \end{tikzfigure}
  }
  \end{minipage}
  \end{adjustbox}
}
}
\column{0.33}
\block{\Huge Implementation}{
We have implemented a number of these isosurface extraction algorithms in an
easy-to-use library. The result is a C/C++ library we call \texttt{libmc}. We
are developing \texttt{libmc} in parallel with our other undergraduate projects
that should benefit from an isosurface extraction library, such as LIDAR
mapping and generalized Voronoi diagram computation.

\texttt{libmc} is still in the early stages of development. Source code for the
library is available on GitHub at the following URL:
\mbox{\url{https://github.com/auntieNeo/libmc}}
}
\newcommand{\examplesize}{6in}
%\block{Example C++ Usage}{
%\lstinputlisting[language=C]{example.c}
%}
\block{Example Output}{
\begin{minipage}[t]{0.5\linewidth}
\setcounter{figurecounter}{0}
\begin{tikzfigure}[Original Marching Cubes]
\includegraphics[width=\examplesize,height=\examplesize]{../build/MC_SIMPLE_MARCHING_CUBES_cube_cd.png}
\end{tikzfigure}
\setcounter{figurecounter}{2}
\begin{tikzfigure}[Cuberille]
\includegraphics[width=\examplesize,height=\examplesize]{../build/MC_CUBERILLE_cube_cd.png}
\end{tikzfigure}
\end{minipage}
\begin{adjustbox}{valign=t}
\begin{minipage}[t]{0.5\linewidth}
\setcounter{figurecounter}{1}
\begin{tikzfigure}[Marching Cubes Patch]
\includegraphics[width=\examplesize,height=\examplesize]{../build/MC_PATCH_MARCHING_CUBES_cube_cd.png}
\end{tikzfigure}
\setcounter{figurecounter}{3}
\begin{tikzfigure}[Nielson's Dual Marching Cubes]
\includegraphics[width=\examplesize,height=\examplesize]{../build/MC_NIELSON_DUAL_cube_cd.png}
\end{tikzfigure}
\end{minipage}
\end{adjustbox}
}
\block{Contact}{
I can be reached at \mbox{\href{mailto:glinjona@isu.edu}{glinjona@isu.edu}}. I also
have a personal web page at \mbox{\url{http://glines.net}}.

My research advisor, Dr. John Edwards, keeps a faculty page at the URL
\mbox{\url{http://www2.cose.isu.edu/~edwajohn/}} and can be contacted via email at
\mbox{\href{mailto:edwajohn@isu.edu}{edwajohn@isu.edu}}.
}
%\begin{subcolumns}
%\subcolumn{0.5}
%\block{Generalizing the Transvoxel Algorithm}{
%The Transvoxel algorithm, as described by Lengyel in FIXME,
%\\ \\
%\begin{tikzpicture}[scale=3]
%  \draw (0, 0) rectangle (1, 1);
%  \draw (1, 0) rectangle (2, 1);
%  \draw (0, 1) rectangle (1, 2);
%  \draw (1, 1) rectangle (2, 2);
%
%  \draw (2, 2) rectangle (3, 3);
%  \draw (3, 2) rectangle (4, 3);
%  \draw (2, 3) rectangle (3, 4);
%  \draw (3, 3) rectangle (4, 4);
%
%  \draw (0, 2) rectangle (1, 3);
%  \draw (1, 2) rectangle (2, 3);
%  \draw (0, 3) rectangle (1, 4);
%  \draw (1, 3) rectangle (2, 4);
%
%  \draw (2.5, 0) rectangle (4, 1.5);
%
%  \begin{scope}[ultra thick]
%    \filldraw[fill=green!20!white, draw=green!50!black] (2, 2) -- (2, 0) -- (2.5, 0) -- (2.5, 1.5) -- (2, 2);
%    \filldraw[fill=green!20!white, draw=green!50!black] (2, 2) -- (4, 2) -- (4, 1.5) -- (2.5, 1.5) -- (2, 2) node [below=1.6em, right=3.7em] {Transition Cell};
%  \end{scope}
%\end{tikzpicture} \\ \\
%
%\begin{tikzpicture}[scale=1.2]
%  % Draw a border around the quadtree figure
%  \begin{scope}[ultra thick]
%    \draw (0, 0) rectangle (8, 12);
%  \end{scope}
%  % Draw the quadtree nodes at three levels
%  \begin{scope}[thick]
%    \draw (0, 0) rectangle (4, 4);
%    \draw (4, 0) rectangle (8, 4);
%    \draw (0, 4) rectangle (4, 8);
%    \draw (4, 4) rectangle (8, 8);
%    \draw (0, 8) rectangle (4, 12);
%    \draw (4, 8) rectangle (8, 12);
%  \end{scope}
%  \begin{scope}
%    \draw (0, 0) rectangle (2, 2);
%    \draw (2, 0) rectangle (4, 2);
%    \draw (0, 2) rectangle (2, 4);
%    \draw (2, 2) rectangle (4, 4);
%
%    \draw (0, 4) rectangle (2, 6);
%    \draw (2, 4) rectangle (4, 6);
%    \draw (0, 6) rectangle (2, 8);
%    \draw (2, 6) rectangle (4, 8);
%  \end{scope}
%  \begin{scope}[very thin]
%    \draw (2, 6) rectangle (3, 7);
%    \draw (3, 6) rectangle (4, 7);
%    \draw (2, 7) rectangle (3, 8);
%    \draw (3, 7) rectangle (4, 8);
%  \end{scope}
%  % Draw the transition cells, being careful to draw the larger ones first
%  \begin{scope}[fill=green!20!white, draw=green!50!black, thick]
%    % Large transitions
%    \filldraw (4, 0) -- (6, 1) -- (6, 4) -- (4, 4) -- (4, 0);
%    \filldraw (4, 4) -- (6, 4) -- (6, 7) -- (4, 8) -- (4, 4);
%    \filldraw (4, 8) -- (3, 10) -- (0, 10) -- (0, 8) -- (4, 8);
%    % Small transitions
%    \filldraw (4, 4) -- (5, 5) -- (5, 6) -- (4, 6) -- (4, 4);
%    \filldraw (4, 6) -- (5, 6) -- (5, 7) -- (4, 8) -- (4, 6);
%    \filldraw (4, 8) -- (3, 9) -- (2, 9) -- (2, 8) -- (4, 8);
%    \filldraw (0, 8) -- (2, 8) -- (2, 9) -- (0, 9) -- (0, 8);
%    % Smallest transitions
%    \filldraw (2, 6) -- (2.5, 5.5) -- (3.5, 5.5) -- (4, 6) -- (2, 6);
%    \filldraw (2, 6) -- (1.5, 6.5) -- (1.5, 7.5) -- (2, 8) -- (2, 6);
%  \end{scope}
%\end{tikzpicture}
%}
%\subcolumn{0.5}
%%\begin{minipage}[t]{0.1666\linewidth}
%%\begin{tikzfigure}
%%\includegraphics[width=1in,height=1in]{../build/MC_TRANSVOXEL_transition_cell_0x001.png}
%%\end{tikzfigure}
%%\end{minipage}
%%\begin{adjustbox}{valign=t}
%%\begin{minipage}[t]{0.5\linewidth}
%%\begin{tikzfigure}[Cuberille]
%%\includegraphics[width=4in,height=4in]{../build/MC_CUBERILLE_cube_cd.png}
%%\end{tikzfigure}
%%\begin{tikzfigure}[Nielson's Dual Marching Cubes]
%%\includegraphics[width=4in,height=4in]{../build/MC_NIELSON_DUAL_cube_cd.png}
%%\end{tikzfigure}
%%\end{minipage}
%%\end{adjustbox}
%\end{subcolumns}
\end{columns}
\end{document}
